\section{微分方程式}\label{節:微分方程式}
これまでは,関数が与えられたときに導関数を求めた.逆に,導関数についての条件から関数を決めることもできる.例えば
\[
 y'=ky
\]
を満たす関数$y$を考えよう.\,$k$は定数である.未知数が数ではなく関数であり,その導関数も含む方程式を\textbf{微分方程式}という.

$y=\e^{kx}$ならこの式を満たす.定数倍した$C\e^{kx}$も満たすが,ほかの形の解はないだろうか.そのことを調べるため,$y$を既に見つけた解$\e^{kx}$で割ってみよう.\,$y=C\e^{kx}$なら,この商は定数$C$になるはずである.解$y$に対して実際に微分すると,
\[
 \frac{\dd}{\dd x}(\e^{-kx}y)
 =\e^{-kx}(y'-ky)=0
\]
となる.平均値の定理から,導関数が$0$の関数は一つの区間で定数である.したがって$\e^{-kx}y=C$より$y=C\e^{kx}$となる.\,$y(0)=A$も与えられれば$C=A$であり,解は一つに定まる.

等しい間隔$h$での値だけを取り出して$a_n=y(nh)$とおけば,
\[
 a_{n+1}=\e^{kh}a_n
\]
である.この条件から見れば等比数列となっている.微分方程式を学ぶ前から知っていた等比的な増加も,間の値を指数関数で表すことができる.

もう一つ,\,$y''=-y$を考えてみよう.\,$\sin x,\cos x$はいずれも二回微分すると符号が変わるので,
\[
 y=A\cos x+B\sin x
\]
は解である.さらに$y(0)=A$,\,$y'(0)=B$を満たす.これ以外の解がないことも確かめられる.同じ初期条件をもつ別の解との差を$z$とすると,
\[
 z''=-z,\qquad z(0)=z'(0)=0
\]
である.\,$z''=-z$の両辺に$z'$を掛けると,$z'z''=-zz'$となる.左辺は$\frac12(z')^2$の微分,右辺は$-\frac12z^2$の微分だから,二つを同じ側に集めた$z^2+(z')^2$を調べればよい.実際,
\[
 \{z^2+(z')^2\}'=2zz'+2z'z''=2z'(z+z'')=0.
\]
したがってこの式は定数で,\,$x=0$で$0$だから常に$0$となる.二つの平方の和が$0$なので$z=0$である.微分方程式の一般的な理論を使わずに,解の形と一意性を確かめられた.

時刻$t$での位置を$s(t)$とすると,速度は$s'(t)$であり,速度の変化率$s''(t)$を\textbf{加速度}という.例えば$s''(t)=-s(t)$は,原点から離れるほど大きく,原点へ向かう加速度が生じる運動を表す.先ほどの結果から,位置は$s(t)=A\cos t+B\sin t$となり,周期的に往復する.このような運動を\textbf{単振動}という.同じ微分方程式は,変数を時刻と考えない場合にも,二回微分したときの関係によって三角関数を特徴付けている.

ここまでは,微分から分かることを調べた.次章では,導入で扱った面積を一般の関数について定め,その変化率が元の関数になることを示す.そこから,微分の逆によって積分を計算する方法を考えよう.
